% 可在 % 可由 % 可由 % 可由 \input{extension.tex} 纳入主稿；要求 amsmath, amssymb, amsthm。
\section{不等指数的共同中心传递}\label{sec:weighted-exponents}

本节把等指数乘积的缺陷坐标推广到任意正指数。该推广仍要求各因子在同一个中心方体上取平均；分别取上确界的多线性模型不适用本节结论。

令 $m\ge2$、$a_1,\ldots,a_m>0$，并在 $0\le s\le1$ 定义
\[
 P_{\boldsymbol a}(s)=
 \max\left\{\prod_{i=1}^m(1-d_i)^{a_i}:
 0\le d_i<1,\ \sum_{i=1}^m d_i=s\right\}.
 \tag{W.1}
\]
函数 $P_{\boldsymbol a}$ 连续、严格递减，且 $P_{\boldsymbol a}(0)=1$。因此当
\[
 P_{\boldsymbol a}(1)<\lambda<1
 \tag{W.2}
\]
时，存在唯一的 $s_{\boldsymbol a}(\lambda)\in(0,1)$ 满足
\[
 P_{\boldsymbol a}(s_{\boldsymbol a}(\lambda))=\lambda.
 \tag{W.3}
\]

\begin{lemma}[水位分配]\label{lem:water-filling}
若 $0<s<1$，存在唯一的 $\rho>0$ 使
\[
 \sum_{i=1}^m\left(1-\frac{a_i}{\rho}\right)_+=s.
 \tag{W.4}
\]
式 \textup{(W.1)} 的唯一最大点为
\[
 d_i^*(s)=\left(1-\frac{a_i}{\rho}\right)_+.
 \tag{W.5}
\]
\end{lemma}

\begin{proof}
在单纯形的相对内部，目标的对数为严格凹函数
\(\sum_i a_i\log(1-d_i)\)。KKT 条件给出：活跃坐标满足
\(a_i/(1-d_i)=\rho\)，非活跃坐标满足 $a_i\ge\rho$。于是得到
\textup{(W.4)}--\textup{(W.5)}。左边作为 $\rho$ 的函数连续递增，从 $0$ 取到大于 $1$ 的值，因而 $\rho$ 唯一。严格递减性和连续性也随之得到：增加总缺陷后，把一个正缺陷坐标减小即可严格增大乘积。
\end{proof}

设 $\mu$ 是非原子、局部有限的 Radon 测度，零质量方体一律略去，
\[
 \mathcal M_{c,\mu}^{\boldsymbol a}(f_1,\ldots,f_m)(x)
 =\sup_{r>0}\prod_{i=1}^m
 \left(\frac{1}{\mu(Q(x,r))}\int_{Q(x,r)}|f_i|\,d\mu\right)^{a_i}.
\]
定义共同中心常数
\[
 C_{c,\mu}^{\boldsymbol a}(\lambda)=
 \sup_{\substack{E_1,\ldots,E_m\ \mathrm{measurable}\\
 F=\cap_iE_i,\ 0<\mu(F)<\infty}}
 \frac{\mu\{\mathcal M_{c,\mu}^{\boldsymbol a}
 (\mathbf1_{E_1},\ldots,\mathbf1_{E_m})>\lambda\}}{\mu(F)}.
\]

\begin{theorem}[不等指数共同中心精确传递]\label{thm:weighted-transfer}
在 \textup{(W.2)} 的区间内，
\[
 C_{c,\mu}^{\boldsymbol a}(\lambda)=
 C_{c,\mu}^{\mathrm{int}}\bigl(1-s_{\boldsymbol a}(\lambda)\bigr).
 \tag{W.6}
\]
若 $d\mu=w\,dx$，上式适用于任意非零非负局部可积权。对外层模型，把左侧平均改为 Lebesgue 平均、光晕测量改为 $w\,dx$，并要求 $0<|F|<\infty$，仍有同一等式。
\end{theorem}

\begin{proof}
写 $s=s_{\boldsymbol a}(\lambda)$。若在共同见证方体 $Q$ 上
\[
 p_i=\frac{\mu(E_i\cap Q)}{\mu(Q)},\qquad d_i=1-p_i,
\]
且 $\prod_i p_i^{a_i}>\lambda$，若总缺陷至少为 $s$，可逐坐标减小缺陷到总和恰为 $s$；乘积只能增大，且至多为 $P_{\boldsymbol a}(s)=\lambda$，矛盾。因此定义 \textup{(W.1)} 给出
\(\sum_i d_i<s\)。又 $Q\setminus F\subset\bigcup_i(Q\setminus E_i)$，故
\[
 \frac{\mu(F\cap Q)}{\mu(Q)}
 \ge1-\sum_i d_i>1-s.
\]
这给出 \textup{(W.6)} 的“$\le$”方向。

反之，令 $K$ 为严格线性光晕
\(\{M_{c,\mu}\mathbf1_F>1-s\}$ 的紧子集。我们使用如下有限稳定事实：可选有限个参考方体以及正余量，使 $K$ 被相应的小中心邻域覆盖；每个邻域中的中心方体仍有 $F$-缺陷严格小于 $s$；并且参考方体及这些小中心扰动的对称差的 $\mu$-质量小到小于预先指定的余量。先把严格见证半径略微增大，避开至多可数个具有正质量边界的半径，使每个参考方体边界的 $\mu$-质量为零。严格性在充分小的增大下保持。此时这一事实由严格不等式、局部有限性、从上连续性和
\(\mu(Q(x,r)\mathbin\triangle Q(y,r))\to0\) 得到。

取包含全部微扰方体的有界盒 $D$。把有限参考方体所生成的 Boolean 原子上的 $D\setminus F$ 按下列比例切分：
\[
 \frac{d_1^*(s)}s,\ldots,\frac{d_m^*(s)}s.
\]
非原子性保证每个原子可作这种有限分割；在 $D$ 外把 $F^c$ 任意补入这些部分，得到 $F^c=A_1\mathbin{\dot\cup}\cdots\mathbin{\dot\cup}A_m$，并置 $E_i=\mathbb R^n\setminus A_i$。对未扰动的参考方体，
\[
 \frac{\mu(Q\setminus E_i)}{\mu(Q)}=
 \frac{d_i^*(s)}s\frac{\mu(Q\setminus F)}{\mu(Q)}<d_i^*(s)\quad(d_i^*(s)>0).
 \tag{W.7}
\]
有限稳定性中的余量使 \textup{(W.7)} 在所覆盖的全部中心扰动方体中仍严格成立。这里把 $D$ 取为包含所有微扰方体的一个有界盒；在该盒内按参考方体的 Boolean 原子分割。若 $d_i^*(s)=0$，令 $A_i\cap D=\varnothing$，故该坐标没有误差。于是每个相关方体满足
\[
 \prod_i\left(\frac{\mu(E_i\cap Q)}{\mu(Q)}\right)^{a_i}
 >\prod_i(1-d_i^*(s))^{a_i}=\lambda.
\]
因此 $K$ 包含于所构造的多线性光晕。由 Radon 内正则性令 $K$ 穷尽线性光晕，得到反向不等式。

外层版本完全相同：对 Lebesgue 测度作有限 Boolean 原子切分，再以 $w\,dx$ 的 Radon 内正则性穷尽光晕。
\end{proof}

\begin{corollary}[等指数特例]
当 $a_1=\cdots=a_m=1$ 时，
\[
 P_{\boldsymbol a}(s)=(1-s/m)^m,\qquad
 s_{\boldsymbol a}(\lambda)=m(1-\lambda^{1/m}),
\]
从而恢复等指数的共同中心缺陷传递。
\end{corollary}

\begin{remark}[外层临界发散]
若 $m\ge2$，则对每个非零局部可积权 $w$ 和
\(0<\lambda\le P_{\boldsymbol a}(1)\)，外层常数为无穷。证明取沿第一坐标的周期分割，其各部分的比例为 $d_i^*(1)$；把一个有正 $w$-质量的有限集 $F$ 放进某个正比例部分，并令 $E_i=(\mathbb R^n\setminus A_i)\cup F$。每个充分大的整数边长中心方体给出严格大于 $P_{\boldsymbol a}(1)$ 的共同见证，故光晕为全空间。有限总质量情形再令 $w(F)\downarrow0$。该论证改变内禀平均分母，不能推出内禀模型的临界发散。
\end{remark}
 纳入主稿；要求 amsmath, amssymb, amsthm。
\section{不等指数的共同中心传递}\label{sec:weighted-exponents}

本节把等指数乘积的缺陷坐标推广到任意正指数。该推广仍要求各因子在同一个中心方体上取平均；分别取上确界的多线性模型不适用本节结论。

令 $m\ge2$、$a_1,\ldots,a_m>0$，并在 $0\le s\le1$ 定义
\[
 P_{\boldsymbol a}(s)=
 \max\left\{\prod_{i=1}^m(1-d_i)^{a_i}:
 0\le d_i<1,\ \sum_{i=1}^m d_i=s\right\}.
 \tag{W.1}
\]
函数 $P_{\boldsymbol a}$ 连续、严格递减，且 $P_{\boldsymbol a}(0)=1$。因此当
\[
 P_{\boldsymbol a}(1)<\lambda<1
 \tag{W.2}
\]
时，存在唯一的 $s_{\boldsymbol a}(\lambda)\in(0,1)$ 满足
\[
 P_{\boldsymbol a}(s_{\boldsymbol a}(\lambda))=\lambda.
 \tag{W.3}
\]

\begin{lemma}[水位分配]\label{lem:water-filling}
若 $0<s<1$，存在唯一的 $\rho>0$ 使
\[
 \sum_{i=1}^m\left(1-\frac{a_i}{\rho}\right)_+=s.
 \tag{W.4}
\]
式 \textup{(W.1)} 的唯一最大点为
\[
 d_i^*(s)=\left(1-\frac{a_i}{\rho}\right)_+.
 \tag{W.5}
\]
\end{lemma}

\begin{proof}
在单纯形的相对内部，目标的对数为严格凹函数
\(\sum_i a_i\log(1-d_i)\)。KKT 条件给出：活跃坐标满足
\(a_i/(1-d_i)=\rho\)，非活跃坐标满足 $a_i\ge\rho$。于是得到
\textup{(W.4)}--\textup{(W.5)}。左边作为 $\rho$ 的函数连续递增，从 $0$ 取到大于 $1$ 的值，因而 $\rho$ 唯一。严格递减性和连续性也随之得到：增加总缺陷后，把一个正缺陷坐标减小即可严格增大乘积。
\end{proof}

设 $\mu$ 是非原子、局部有限的 Radon 测度，零质量方体一律略去，
\[
 \mathcal M_{c,\mu}^{\boldsymbol a}(f_1,\ldots,f_m)(x)
 =\sup_{r>0}\prod_{i=1}^m
 \left(\frac{1}{\mu(Q(x,r))}\int_{Q(x,r)}|f_i|\,d\mu\right)^{a_i}.
\]
定义共同中心常数
\[
 C_{c,\mu}^{\boldsymbol a}(\lambda)=
 \sup_{\substack{E_1,\ldots,E_m\ \mathrm{measurable}\\
 F=\cap_iE_i,\ 0<\mu(F)<\infty}}
 \frac{\mu\{\mathcal M_{c,\mu}^{\boldsymbol a}
 (\mathbf1_{E_1},\ldots,\mathbf1_{E_m})>\lambda\}}{\mu(F)}.
\]

\begin{theorem}[不等指数共同中心精确传递]\label{thm:weighted-transfer}
在 \textup{(W.2)} 的区间内，
\[
 C_{c,\mu}^{\boldsymbol a}(\lambda)=
 C_{c,\mu}^{\mathrm{int}}\bigl(1-s_{\boldsymbol a}(\lambda)\bigr).
 \tag{W.6}
\]
若 $d\mu=w\,dx$，上式适用于任意非零非负局部可积权。对外层模型，把左侧平均改为 Lebesgue 平均、光晕测量改为 $w\,dx$，并要求 $0<|F|<\infty$，仍有同一等式。
\end{theorem}

\begin{proof}
写 $s=s_{\boldsymbol a}(\lambda)$。若在共同见证方体 $Q$ 上
\[
 p_i=\frac{\mu(E_i\cap Q)}{\mu(Q)},\qquad d_i=1-p_i,
\]
且 $\prod_i p_i^{a_i}>\lambda$，若总缺陷至少为 $s$，可逐坐标减小缺陷到总和恰为 $s$；乘积只能增大，且至多为 $P_{\boldsymbol a}(s)=\lambda$，矛盾。因此定义 \textup{(W.1)} 给出
\(\sum_i d_i<s\)。又 $Q\setminus F\subset\bigcup_i(Q\setminus E_i)$，故
\[
 \frac{\mu(F\cap Q)}{\mu(Q)}
 \ge1-\sum_i d_i>1-s.
\]
这给出 \textup{(W.6)} 的“$\le$”方向。

反之，令 $K$ 为严格线性光晕
\(\{M_{c,\mu}\mathbf1_F>1-s\}$ 的紧子集。我们使用如下有限稳定事实：可选有限个参考方体以及正余量，使 $K$ 被相应的小中心邻域覆盖；每个邻域中的中心方体仍有 $F$-缺陷严格小于 $s$；并且参考方体及这些小中心扰动的对称差的 $\mu$-质量小到小于预先指定的余量。先把严格见证半径略微增大，避开至多可数个具有正质量边界的半径，使每个参考方体边界的 $\mu$-质量为零。严格性在充分小的增大下保持。此时这一事实由严格不等式、局部有限性、从上连续性和
\(\mu(Q(x,r)\mathbin\triangle Q(y,r))\to0\) 得到。

取包含全部微扰方体的有界盒 $D$。把有限参考方体所生成的 Boolean 原子上的 $D\setminus F$ 按下列比例切分：
\[
 \frac{d_1^*(s)}s,\ldots,\frac{d_m^*(s)}s.
\]
非原子性保证每个原子可作这种有限分割；在 $D$ 外把 $F^c$ 任意补入这些部分，得到 $F^c=A_1\mathbin{\dot\cup}\cdots\mathbin{\dot\cup}A_m$，并置 $E_i=\mathbb R^n\setminus A_i$。对未扰动的参考方体，
\[
 \frac{\mu(Q\setminus E_i)}{\mu(Q)}=
 \frac{d_i^*(s)}s\frac{\mu(Q\setminus F)}{\mu(Q)}<d_i^*(s)\quad(d_i^*(s)>0).
 \tag{W.7}
\]
有限稳定性中的余量使 \textup{(W.7)} 在所覆盖的全部中心扰动方体中仍严格成立。这里把 $D$ 取为包含所有微扰方体的一个有界盒；在该盒内按参考方体的 Boolean 原子分割。若 $d_i^*(s)=0$，令 $A_i\cap D=\varnothing$，故该坐标没有误差。于是每个相关方体满足
\[
 \prod_i\left(\frac{\mu(E_i\cap Q)}{\mu(Q)}\right)^{a_i}
 >\prod_i(1-d_i^*(s))^{a_i}=\lambda.
\]
因此 $K$ 包含于所构造的多线性光晕。由 Radon 内正则性令 $K$ 穷尽线性光晕，得到反向不等式。

外层版本完全相同：对 Lebesgue 测度作有限 Boolean 原子切分，再以 $w\,dx$ 的 Radon 内正则性穷尽光晕。
\end{proof}

\begin{corollary}[等指数特例]
当 $a_1=\cdots=a_m=1$ 时，
\[
 P_{\boldsymbol a}(s)=(1-s/m)^m,\qquad
 s_{\boldsymbol a}(\lambda)=m(1-\lambda^{1/m}),
\]
从而恢复等指数的共同中心缺陷传递。
\end{corollary}

\begin{remark}[外层临界发散]
若 $m\ge2$，则对每个非零局部可积权 $w$ 和
\(0<\lambda\le P_{\boldsymbol a}(1)\)，外层常数为无穷。证明取沿第一坐标的周期分割，其各部分的比例为 $d_i^*(1)$；把一个有正 $w$-质量的有限集 $F$ 放进某个正比例部分，并令 $E_i=(\mathbb R^n\setminus A_i)\cup F$。每个充分大的整数边长中心方体给出严格大于 $P_{\boldsymbol a}(1)$ 的共同见证，故光晕为全空间。有限总质量情形再令 $w(F)\downarrow0$。该论证改变内禀平均分母，不能推出内禀模型的临界发散。
\end{remark}
 纳入主稿；要求 amsmath, amssymb, amsthm。
\section{不等指数的共同中心传递}\label{sec:weighted-exponents}

本节把等指数乘积的缺陷坐标推广到任意正指数。该推广仍要求各因子在同一个中心方体上取平均；分别取上确界的多线性模型不适用本节结论。

令 $m\ge2$、$a_1,\ldots,a_m>0$，并在 $0\le s\le1$ 定义
\[
 P_{\boldsymbol a}(s)=
 \max\left\{\prod_{i=1}^m(1-d_i)^{a_i}:
 0\le d_i<1,\ \sum_{i=1}^m d_i=s\right\}.
 \tag{W.1}
\]
函数 $P_{\boldsymbol a}$ 连续、严格递减，且 $P_{\boldsymbol a}(0)=1$。因此当
\[
 P_{\boldsymbol a}(1)<\lambda<1
 \tag{W.2}
\]
时，存在唯一的 $s_{\boldsymbol a}(\lambda)\in(0,1)$ 满足
\[
 P_{\boldsymbol a}(s_{\boldsymbol a}(\lambda))=\lambda.
 \tag{W.3}
\]

\begin{lemma}[水位分配]\label{lem:water-filling}
若 $0<s<1$，存在唯一的 $\rho>0$ 使
\[
 \sum_{i=1}^m\left(1-\frac{a_i}{\rho}\right)_+=s.
 \tag{W.4}
\]
式 \textup{(W.1)} 的唯一最大点为
\[
 d_i^*(s)=\left(1-\frac{a_i}{\rho}\right)_+.
 \tag{W.5}
\]
\end{lemma}

\begin{proof}
在单纯形的相对内部，目标的对数为严格凹函数
\(\sum_i a_i\log(1-d_i)\)。KKT 条件给出：活跃坐标满足
\(a_i/(1-d_i)=\rho\)，非活跃坐标满足 $a_i\ge\rho$。于是得到
\textup{(W.4)}--\textup{(W.5)}。左边作为 $\rho$ 的函数连续递增，从 $0$ 取到大于 $1$ 的值，因而 $\rho$ 唯一。严格递减性和连续性也随之得到：增加总缺陷后，把一个正缺陷坐标减小即可严格增大乘积。
\end{proof}

设 $\mu$ 是非原子、局部有限的 Radon 测度，零质量方体一律略去，
\[
 \mathcal M_{c,\mu}^{\boldsymbol a}(f_1,\ldots,f_m)(x)
 =\sup_{r>0}\prod_{i=1}^m
 \left(\frac{1}{\mu(Q(x,r))}\int_{Q(x,r)}|f_i|\,d\mu\right)^{a_i}.
\]
定义共同中心常数
\[
 C_{c,\mu}^{\boldsymbol a}(\lambda)=
 \sup_{\substack{E_1,\ldots,E_m\ \mathrm{measurable}\\
 F=\cap_iE_i,\ 0<\mu(F)<\infty}}
 \frac{\mu\{\mathcal M_{c,\mu}^{\boldsymbol a}
 (\mathbf1_{E_1},\ldots,\mathbf1_{E_m})>\lambda\}}{\mu(F)}.
\]

\begin{theorem}[不等指数共同中心精确传递]\label{thm:weighted-transfer}
在 \textup{(W.2)} 的区间内，
\[
 C_{c,\mu}^{\boldsymbol a}(\lambda)=
 C_{c,\mu}^{\mathrm{int}}\bigl(1-s_{\boldsymbol a}(\lambda)\bigr).
 \tag{W.6}
\]
若 $d\mu=w\,dx$，上式适用于任意非零非负局部可积权。对外层模型，把左侧平均改为 Lebesgue 平均、光晕测量改为 $w\,dx$，并要求 $0<|F|<\infty$，仍有同一等式。
\end{theorem}

\begin{proof}
写 $s=s_{\boldsymbol a}(\lambda)$。若在共同见证方体 $Q$ 上
\[
 p_i=\frac{\mu(E_i\cap Q)}{\mu(Q)},\qquad d_i=1-p_i,
\]
且 $\prod_i p_i^{a_i}>\lambda$，若总缺陷至少为 $s$，可逐坐标减小缺陷到总和恰为 $s$；乘积只能增大，且至多为 $P_{\boldsymbol a}(s)=\lambda$，矛盾。因此定义 \textup{(W.1)} 给出
\(\sum_i d_i<s\)。又 $Q\setminus F\subset\bigcup_i(Q\setminus E_i)$，故
\[
 \frac{\mu(F\cap Q)}{\mu(Q)}
 \ge1-\sum_i d_i>1-s.
\]
这给出 \textup{(W.6)} 的“$\le$”方向。

反之，令 $K$ 为严格线性光晕
\(\{M_{c,\mu}\mathbf1_F>1-s\}$ 的紧子集。我们使用如下有限稳定事实：可选有限个参考方体以及正余量，使 $K$ 被相应的小中心邻域覆盖；每个邻域中的中心方体仍有 $F$-缺陷严格小于 $s$；并且参考方体及这些小中心扰动的对称差的 $\mu$-质量小到小于预先指定的余量。先把严格见证半径略微增大，避开至多可数个具有正质量边界的半径，使每个参考方体边界的 $\mu$-质量为零。严格性在充分小的增大下保持。此时这一事实由严格不等式、局部有限性、从上连续性和
\(\mu(Q(x,r)\mathbin\triangle Q(y,r))\to0\) 得到。

取包含全部微扰方体的有界盒 $D$。把有限参考方体所生成的 Boolean 原子上的 $D\setminus F$ 按下列比例切分：
\[
 \frac{d_1^*(s)}s,\ldots,\frac{d_m^*(s)}s.
\]
非原子性保证每个原子可作这种有限分割；在 $D$ 外把 $F^c$ 任意补入这些部分，得到 $F^c=A_1\mathbin{\dot\cup}\cdots\mathbin{\dot\cup}A_m$，并置 $E_i=\mathbb R^n\setminus A_i$。对未扰动的参考方体，
\[
 \frac{\mu(Q\setminus E_i)}{\mu(Q)}=
 \frac{d_i^*(s)}s\frac{\mu(Q\setminus F)}{\mu(Q)}<d_i^*(s)\quad(d_i^*(s)>0).
 \tag{W.7}
\]
有限稳定性中的余量使 \textup{(W.7)} 在所覆盖的全部中心扰动方体中仍严格成立。这里把 $D$ 取为包含所有微扰方体的一个有界盒；在该盒内按参考方体的 Boolean 原子分割。若 $d_i^*(s)=0$，令 $A_i\cap D=\varnothing$，故该坐标没有误差。于是每个相关方体满足
\[
 \prod_i\left(\frac{\mu(E_i\cap Q)}{\mu(Q)}\right)^{a_i}
 >\prod_i(1-d_i^*(s))^{a_i}=\lambda.
\]
因此 $K$ 包含于所构造的多线性光晕。由 Radon 内正则性令 $K$ 穷尽线性光晕，得到反向不等式。

外层版本完全相同：对 Lebesgue 测度作有限 Boolean 原子切分，再以 $w\,dx$ 的 Radon 内正则性穷尽光晕。
\end{proof}

\begin{corollary}[等指数特例]
当 $a_1=\cdots=a_m=1$ 时，
\[
 P_{\boldsymbol a}(s)=(1-s/m)^m,\qquad
 s_{\boldsymbol a}(\lambda)=m(1-\lambda^{1/m}),
\]
从而恢复等指数的共同中心缺陷传递。
\end{corollary}

\begin{remark}[外层临界发散]
若 $m\ge2$，则对每个非零局部可积权 $w$ 和
\(0<\lambda\le P_{\boldsymbol a}(1)\)，外层常数为无穷。证明取沿第一坐标的周期分割，其各部分的比例为 $d_i^*(1)$；把一个有正 $w$-质量的有限集 $F$ 放进某个正比例部分，并令 $E_i=(\mathbb R^n\setminus A_i)\cup F$。每个充分大的整数边长中心方体给出严格大于 $P_{\boldsymbol a}(1)$ 的共同见证，故光晕为全空间。有限总质量情形再令 $w(F)\downarrow0$。该论证改变内禀平均分母，不能推出内禀模型的临界发散。
\end{remark}
 之后纳入主稿；依赖 amsmath, amssymb, amsthm。
\begin{proposition}[有限总质量时的内禀临界发散]\label{prop:finite-mass-critical}
令 $m\ge2$、$a_1,\ldots,a_m>0$，并令 $\mu=w\,dx$ 是非零有限
Radon 测度（故 $\mu$ 非原子）。设 $P_{\boldsymbol a}$ 由
\textup{\rm(W.1)} 定义。对每个
\[
 0<\lambda\le P_{\boldsymbol a}(1)
\]
都有
\[
 C_{c,\mu}^{\boldsymbol a}(\lambda)=\infty.
\]
\end{proposition}

\begin{proof}
记 $M=\mu(\mathbb R^n)\in(0,\infty)$。在闭单纯形
$\{d_i\ge0:\sum_i d_i=1\}$ 上把目标连续延拓为
$\prod_i(1-d_i)^{a_i}$。取 $d_i=1/m$ 可知最大值为正；而边界上
若某个 $d_i=1$，该延拓值为零。因此存在最大点
$\boldsymbol d^*=(d_1^*,\ldots,d_m^*)$，满足
\[
 P_{\boldsymbol a}(1)=\prod_i(1-d_i^*)^{a_i}>0,
 \qquad 0\le d_i^*<1,\qquad \sum_i d_i^*=1.
\]
选取 $d_j^*>0$。非原子性给出 $\mathbb R^n=A_1\mathbin{\dot\cup}
\cdots\mathbin{\dot\cup}A_m$，其中 $\mu(A_i)=d_i^*M$（零比例部分
可取零测度）。任取 $0<t<d_j^*M$，再取 $F\subset A_j$ 使
$\mu(F)=t$，并置
\[
 E_i=A_i^c\cup F\quad(1\le i\le m).
\]
这时 $\bigcap_iE_i=F$，且
\[
 \prod_i\left(\frac{\mu(E_i)}M\right)^{a_i}
 =\left(1-d_j^*+\frac tM\right)^{a_j}
   \prod_{i\ne j}(1-d_i^*)^{a_i}
 >P_{\boldsymbol a}(1)\ge\lambda. \tag{FM.1}
\]
对每个 $x\in\mathbb R^n$，有 $Q(x,R)\uparrow\mathbb R^n$ 以及
$\mu(Q(x,R))\to M>0$。故左边的方体平均乘积收敛到 \textup{(FM.1)}
的左边；对充分大的（依赖于 $x$ 的）$R$，其严格大于 $\lambda$。
所以相应多线性光晕就是 $\mathbb R^n$，并且
\[
 \frac{\mu\{\mathcal M_{c,\mu}^{\boldsymbol a}
 (\mathbf1_{E_1},\ldots,\mathbf1_{E_m})>\lambda\}}{\mu(F)}
 =\frac Mt.
\]
令 $t\downarrow0$ 即得结论。
\end{proof}
